\section{Threats to Validity and Research Ethics}
\label{sec:threats_to_validity}
We organize limitations around construct, internal, external, and conclusion validity, informed by empirical reporting guidance~\cite{ralph2020}.

\subsection{Construct Validity}
Registry maintainer arrays supply username associations, not historical publication authority, token scope, account independence, or a peer-review policy. One maintainer does not prove absence of human review; several do not prove enforced review. Uniform and heterogeneous source weights are modeling assumptions. The perfect hardening model removes selected account origins, not package nodes or paths from other origins.

The workflow observation is a directory lookup at HEAD. We did not retain workflow YAML, job execution, permission scopes, historical release links, or consumer lockfiles. Dependency and workflow association therefore establishes neither installation nor vulnerable-code execution or secret access. CI dependency studies~\cite{weeraddana2024} and finer-grained vulnerability analysis~\cite{gao2025emse} motivate those distinctions without supplying execution evidence for our endpoints.

The main exact-version graph is an observed deps.dev edge union, while the temporal companion is a bounded highest-satisfying SemVer approximation. They are different representations. Neither reconstructs npm hoisting, installed versions, overrides, or complete historical dependency trees. Runtime vulnerabilities such as the Node.js Brotli disclosure~\cite{nodejsbrotli2026} are outside the package graph. Residual thresholds are hypothetical source-count rules, without signer groups, role constraints, or measured enforcement.

\subsection{Internal and Observational Validity}
Independent traversals, transposed source incidence, small-graph cases, and generated result checks validate computations on retained evidence. They cannot recover omitted observations. SCC condensation preserves reach; it is not a source of measured cycle-removal benefit. Ties use ascending username order, not portfolio size, except where portfolio size defines the selection rule itself.

Repository normalization is lexical. Excluding inconsistent positive observations tests sensitivity; it does not validate the other remote mappings or repair historical flags. Both main representations mix contexts and acquisition dates. Unknown workflow states and ambiguous zeros can omit eligible endpoints. Historical availability counts describe retained records, including a repeated tuple, and do not establish independent incidents or verified release authority.

\subsection{External and Conclusion Validity}
The high-impact-seeded expansion is neither registry-wide nor a random sample of npm installations. Public workflow-associated repositories cannot represent private consumer infrastructure without further evidence. Results describe this graph and explicit scenarios. They supply no measured compromise probabilities or causal control effects.

Random orders quantify selection variability. Orders are reused across nested budgets and endpoint definitions, so the resulting 6,000 rows are not independent ecosystem samples. Monte Carlo errors describe simulation means, not npm-wide uncertainty. The seven scalarization settings are exploratory tradeoff points, without a certified global Pareto frontier or an empirically estimated stakeholder preference. Heterogeneous weights are structural sensitivities, not fitted risk priors. Claims about efficacy require a separate intervention study.

\subsection{Ethics and Reporting}
The analyses read retained public metadata and execute no package payloads, credential probes, or exploitation. Prospective captures omit email addresses and package bodies. Public usernames illustrate structural overlap without implying present-day insecurity, negligence, or historical publishing authority. Incident associations do not validate prediction. Artifact checksums identify retained bytes rather than certify authenticity. Assistance is disclosed, and authors must review the claims, sources, code, and applicable submission requirements.
